\documentclass[11pt]{article}
\usepackage[margin=1in]{geometry}
\usepackage{amsmath,amssymb,amsthm}
\usepackage{booktabs}
\usepackage{bm}

% --- Macros, consistent with meta_simulator.tex / notations.tex ---
\newcommand{\rmB}{\mathbf{B}}
\newcommand{\rmX}{\mathbf{X}}
\newcommand{\rmM}{\mathbf{M}}
\newcommand{\rmY}{\mathbf{Y}}
\newcommand{\rmT}{\mathbf{T}}
\newcommand{\rvs}{\bm{s}}
\newcommand{\rvtheta}{\bm{\theta}}
\providecommand{\cref}[1]{Eq.~(\ref{#1})}
\providecommand{\autoref}[1]{\ref{#1}}

\newtheorem{proposition}{Proposition}
\newtheorem{remark}{Remark}

\title{Amortized Theory Discovery\\
\large Structure inference over design configurations with CogFormer}
\author{CogFormer-2 (working notes)}
\date{}

\begin{document}
\maketitle

\begin{abstract}
\noindent
CogFormer amortizes the posterior $p(\rmB \mid \rmY, \rmX, \rmM, \mathcal{F})$ over
coefficients \emph{given} a design configuration $\rmM$. This note formalizes the
complementary problem: given real data, infer $\rmM$ itself. We treat a design
configuration as the numerical encoding of a \emph{theory}, cast theory discovery as
Bayesian structure inference over a combinatorial space of $2^{\,\mathcal{O}(RD_\mathcal{F})}$
hypotheses, and show that the same meta-simulator that trains CogFormer also supplies
free supervision for an amortized posterior over that space. The result is a posterior
over theories obtainable in a single forward pass, from which Bayes factors, posterior
inclusion probabilities, and model-averaged parameter posteriors follow. We also
formalize the limit of the construction: the search is closed over a specified model
space, so an explicit \emph{theory gap} statistic is required to detect that no theory
in the space is adequate.
\end{abstract}

%==============================================================================
\section{Setup}
%==============================================================================

We inherit the generative meta-model of \cref{eq:generative} and the meta-simulator of
\cref{eq:meta-simulator}. A single dataset is generated as
%
\begin{equation}
\label{eq:td-generative}
\mathcal{F} \sim p(\mathcal{F} \mid \mathcal{C}),\quad
(\rmM, \rmX) \sim p(\rmM, \rmX \mid \mathcal{F}),\quad
\rmB \sim p(\rmB \mid \rmM, \mathcal{F}),\quad
\rmY \sim p\big(\rmY \mid g(\rmX(\rmM \odot \rmB)), \mathcal{F}\big),
\end{equation}
%
with $\rmM \in \{0,1\}^{R \times D_\mathcal{F}}$ the parameter mask, $\rmX \in \mathbb{R}^{N \times R}$
the design matrix, and $\rmB \in \mathbb{R}^{R \times D_\mathcal{F}}$ the coefficient matrix.
CogFormer targets the conditional $p(\rmB \mid \mathcal{A})$ with conditioning set
$\mathcal{A} = (\rmY, \rmX, \rmM, \mathcal{F})$. Throughout this note we write
%
\begin{equation}
\mathcal{A}_0 = (\rmY, \rmX)
\end{equation}
%
for the \emph{reduced} conditioning set in which the structural objects $\rmM$ and
$\mathcal{F}$ have been withheld. Theory discovery is inference over exactly the
objects that $\mathcal{A}\setminus\mathcal{A}_0$ contains.

%==============================================================================
\section{Theories as block-level structures}
%==============================================================================

A scientific claim is not made about individual design columns. ``The speed--accuracy
instruction affects the decision threshold'' is a statement about a \emph{regressor} and
a \emph{parameter}, whereas a categorical regressor with $k_c$ levels occupies
$C_{\max}-1$ columns of $\rmX$ and an interaction occupies $(C_{\max}-1)^2$. Inference
at the column level would fragment a single claim across a dummy block and would make
the hypothesis space depend on the realized category counts.

We therefore define theories at the level of blocks. Let $c \in \{0,1,\dots,I^\ast\}$
index design blocks, with $c=0$ the intercept, $c = 1,\dots,I$ the main effects, and
$c > I$ the interaction blocks, and let $\Pi$ denote the expansion that maps a block to
its constituent columns of $\rmX$.

\begin{quote}
\textbf{Definition (theory).} A \emph{theory} is a pair
$\mathcal{H} = (\mathcal{F}, \rmT)$ consisting of a model family $\mathcal{F} \in \mathcal{C}$
and a binary \emph{structure matrix}
$\rmT \in \{0,1\}^{(I^\ast+1) \times D_\mathcal{F}}$, where $T_{c,j} = 1$ asserts that
block $c$ enters intrinsic parameter $j$. Its numerical realization is the parameter
mask $\rmM = \Pi(\rmT)$, with $M_{i,j} = T_{c(i),j}$ for the block $c(i)$ containing
column $i$.
\end{quote}

The admissible set $\mathbb{T}_\mathcal{F}$ collects the structures satisfying the
constraints already imposed by the meta-simulator:
%
\begin{align}
\label{eq:td-constraints}
T_{c,j} \le T_{0,j} \quad \forall c > 0
&\qquad \text{(slopes require an active intercept),}\\
T_{0,j} = 1,\;\; T_{c,j} = 0 \;\; \forall c>0, \; j \in \mathcal{J}_{\mathrm{glob}}
&\qquad \text{(global parameters are intercept-only),}\notag\\
T_{c,j} = 0 \;\;\text{if } \Pi(c) \text{ is unrealized in } \rmX
&\qquad \text{(realized-support masking).}\notag
\end{align}
%
Marginalizing over families, the hypothesis space is
$\mathbb{H} = \bigcup_{\mathcal{F} \in \mathcal{C}} \{\mathcal{F}\} \times \mathbb{T}_\mathcal{F}$,
of cardinality $\mathcal{O}\!\left(\sum_\mathcal{F} 2^{(I^\ast+1)D_\mathcal{F}}\right)$.
For a modest design with $I^\ast = 5$ blocks and $D_\mathcal{F} = 6$ intrinsic parameters,
$|\mathbb{T}_\mathcal{F}|$ before constraints is $2^{36} \approx 7\times10^{10}$.

\begin{remark}
The block-level view is what makes the target scientifically interpretable, and it is
also what makes it tractable to report: the quantity a researcher wants is typically a
single entry $T_{c,j}$, not the whole matrix.
\end{remark}

%==============================================================================
\section{The target and its intractability}
%==============================================================================

Given observed data $\mathcal{A}_0^\ast = (\rmY^\ast, \rmX^\ast)$, the object of interest is
the posterior over theories,
%
\begin{equation}
\label{eq:td-target}
p(\mathcal{H} \mid \mathcal{A}_0^\ast)
= \frac{p(\mathcal{F})\,p(\rmT \mid \mathcal{F})\, Z(\mathcal{H})}
       {\sum_{\mathcal{H}' \in \mathbb{H}} p(\mathcal{F}')\,p(\rmT' \mid \mathcal{F}')\, Z(\mathcal{H}')},
\qquad
Z(\mathcal{H}) = \!\int\! p\big(\rmY^\ast \mid \rmB, \rmX^\ast, \Pi(\rmT), \mathcal{F}\big)\,
p(\rmB \mid \rmT, \mathcal{F})\, d\rmB .
\end{equation}
%
The marginal likelihood $Z(\mathcal{H})$ is \emph{doubly} intractable in this setting.
First, for most families in $\mathcal{C}$ the trial-level likelihood
$p(\rmY \mid \Theta, \mathcal{F})$ has no closed form, which is the premise of
simulation-based inference to begin with. Second, even where it does, the integral over
$\rmB$ is a high-dimensional marginalization whose numerical estimation (bridge
sampling, nested sampling, thermodynamic integration) is markedly more expensive than
posterior estimation, and it would have to be repeated for each of the $|\mathbb{H}|$
hypotheses. The enumeration itself is the binding constraint: the per-model cost is
irrelevant when the model count is $10^{10}$.

%==============================================================================
\section{Amortized structure inference}
%==============================================================================

The route around \cref{eq:td-target} is that \cref{eq:td-generative} defines a joint
distribution over $(\mathcal{F}, \rmT, \rmB, \rmX, \rmY)$ that we can sample from
freely, and $p(\mathcal{H} \mid \mathcal{A}_0)$ is simply one of its conditionals. The
mask is not an external specification supplied by the analyst; it is a random variable
of the simulator, drawn hierarchically per simulation. Every training example therefore
already ships with its own structural label, at no additional simulation cost.

Let $q_\phi(\mathcal{H} \mid \mathcal{A}_0)$ be a neural distribution over
$\mathbb{H}$ conditioned on data and design, and define the expected forward
cross-entropy
%
\begin{equation}
\label{eq:td-loss}
\mathcal{L}(\phi) \;=\; -\,\mathbb{E}_{p(\mathcal{F},\rmT,\rmB,\rmX,\rmY)}
\Big[ \log q_\phi\big(\mathcal{F}, \rmT \mid \rmY, \rmX\big) \Big].
\end{equation}

\begin{proposition}[Self-consistency]
\label{prop:td-selfconsistency}
If $q_\phi$ is sufficiently expressive, the minimizer of \cref{eq:td-loss} satisfies
$q_{\phi^\star}(\mathcal{H} \mid \mathcal{A}_0) = p(\mathcal{H} \mid \mathcal{A}_0)$
for $p(\mathcal{A}_0)$-almost all $\mathcal{A}_0$.
\end{proposition}

\begin{proof}[Proof sketch]
Factor the expectation as
$\mathbb{E}_{p(\mathcal{A}_0)}\mathbb{E}_{p(\mathcal{H}\mid\mathcal{A}_0)}[-\log q_\phi]$
and add and subtract the conditional entropy, giving
$\mathcal{L}(\phi) = \mathbb{E}_{p(\mathcal{A}_0)}\big[\mathrm{KL}\big(p(\mathcal{H}\mid\mathcal{A}_0)\,\|\,q_\phi(\mathcal{H}\mid\mathcal{A}_0)\big)\big] + \mathbb{H}[\mathcal{H}\mid\mathcal{A}_0]$.
The second term is independent of $\phi$ and the first is minimized uniquely at zero.
\end{proof}

\Cref{prop:td-selfconsistency} is the standard amortization argument, but its
consequence here is worth stating plainly: the gradient of \cref{eq:td-loss} involves
neither $Z(\mathcal{H})$ nor any likelihood evaluation. The intractable normalizer of
\cref{eq:td-target} is never formed. Marginal likelihoods are learned implicitly, as a
by-product of learning to predict the structure that generated a dataset.

\subsection{Bayes factors and post-hoc prior reweighting}

Because $q_\phi$ approximates a posterior under the simulator's structure prior,
evidence ratios are recoverable by dividing that prior out:
%
\begin{equation}
\label{eq:td-bayesfactor}
\mathrm{BF}_{\mathcal{H}_1 : \mathcal{H}_2}(\mathcal{A}_0^\ast)
= \frac{Z(\mathcal{H}_1)}{Z(\mathcal{H}_2)}
= \frac{q_\phi(\mathcal{H}_1 \mid \mathcal{A}_0^\ast)\,/\,p(\mathcal{H}_1)}
       {q_\phi(\mathcal{H}_2 \mid \mathcal{A}_0^\ast)\,/\,p(\mathcal{H}_2)} .
\end{equation}
%
Two practical consequences follow. First, the mask prior of the meta-simulator
\emph{is} the model prior of the analysis, and must be reported as such. Second,
because \cref{eq:td-bayesfactor} isolates the evidence ratio, an analyst may substitute
a different structure prior $\tilde{p}(\mathcal{H})$ after training by reweighting,
%
\begin{equation}
\label{eq:td-reweight}
\tilde q(\mathcal{H} \mid \mathcal{A}_0^\ast) \;\propto\;
q_\phi(\mathcal{H} \mid \mathcal{A}_0^\ast)\,\frac{\tilde p(\mathcal{H})}{p(\mathcal{H})},
\end{equation}
%
without retraining. Prior sensitivity analysis over model priors thus costs one forward
pass in total rather than one refit per prior.

\begin{remark}[The simulator's priors are already the right ones]
Two choices made for other reasons turn out to be exactly what structure inference
requires. (i) The hierarchical mask prior, in which inclusion rates
$p_0, p_1 \sim \mathrm{Beta}(2,2)$ are resampled per simulation rather than fixed,
is a multiplicity-correcting prior: learning the inclusion rate from the structure
induces the shrinkage in the number of hypotheses that a fixed Bernoulli rate does not
provide. (ii) The variance-budget coefficient prior, under which the realized slope
variance $\sigma^2_{\beta,j} = \sigma_j^2 w_j/(1-w_j)$ is held fixed regardless of the
number of active blocks, makes the coefficient prior scale-invariant to model
dimension. Marginal likelihoods are notoriously sensitive to the coefficient prior
scale, and diffuse independent slope priors penalize larger models through prior spread
alone rather than through fit. The variance budget removes this artifact by
construction.
\end{remark}

\subsection{Parameterizing $q_\phi$}

\paragraph{Factorized.} The simplest choice treats blocks as conditionally
independent given the data,
%
\begin{equation}
\label{eq:td-meanfield}
q_\phi(\rmT \mid \mathcal{A}_0) = \prod_{c,j} \pi_{c,j}^{T_{c,j}} (1-\pi_{c,j})^{1-T_{c,j}},
\qquad \pi_{c,j} = \sigma\!\big(\ell_\phi(c,j;\mathcal{A}_0)\big),
\end{equation}
%
where $\pi_{c,j}$ is the \emph{posterior inclusion probability} (PIP) of block $c$ in
parameter $j$. This costs a single forward pass and directly answers the marginal
question researchers usually ask. Its limitation is structural: it cannot represent
dependence between entries, and dependence is exactly what collinear designs produce.
When two regressors are strongly correlated and either alone explains the data, the
true posterior is bimodal over $\{(1,0),(0,1)\}$; \cref{eq:td-meanfield} reports
$\pi \approx 0.5$ for both and assigns mass to $(1,1)$ and $(0,0)$, which the true
posterior does not support.

\paragraph{Any-order autoregressive.} Retaining the full joint requires an
autoregressive factorization over a permutation $\sigma$ of the $(I^\ast+1)D_\mathcal{F}$
cells,
%
\begin{equation}
\label{eq:td-aoar}
q_\phi(\rmT \mid \mathcal{A}_0) = \prod_{k=1}^{(I^\ast+1)D_\mathcal{F}}
q_\phi\big(T_{\sigma(k)} \mid T_{\sigma(<k)}, \mathcal{A}_0\big),
\end{equation}
%
trained with $\sigma$ drawn uniformly at random per batch so that all conditional
orderings are learned. This yields exact ancestral sampling of structures, exact joint
probabilities for any enumerated hypothesis, and---because the ordering is
arbitrary---the ability to \emph{condition on a partially specified theory}: a
researcher who commits to ``difficulty affects drift rate'' can query the posterior
over everything else by clamping that cell. This last property maps directly onto how
model lattices are actually explored.

\subsection{Cost, and why speed is preserved}

The architecture already carries one decoder token per coefficient cell $(i,j)$, built
from concatenated parameter and regressor positional embeddings and cross-attending to
the encoder summaries $\rmS$; the mask currently enters as a query padding mask over
these tokens. Structure inference reverses the role of that mask: all cells are
presented as active, and each token emits a Bernoulli logit instead of (or alongside) a
velocity. The encoder, which carries the $\mathcal{O}(N)$ cost in the trial dimension,
is evaluated once. The factorized head \cref{eq:td-meanfield} then costs one decoder
pass; the autoregressive head \cref{eq:td-aoar} costs $(I^\ast+1)D_\mathcal{F}$ decoder
passes over a cached encoder state, which for realistic sizes ($\lesssim 100$ cells) is
comparable to a single flow-matching sampling trajectory. No search over $\mathbb{H}$ is
required, because structures are \emph{sampled} rather than enumerated; enumeration is
needed only for reporting the top-$k$ hypotheses.

%==============================================================================
\section{Joint inference and model averaging}
%==============================================================================

Structure and coefficients compose. Writing $q_\psi$ for the trained CogFormer
posterior, the joint posterior over theory and parameters factorizes as
%
\begin{equation}
\label{eq:td-joint}
q(\rmB, \mathcal{H} \mid \mathcal{A}_0^\ast)
= \underbrace{q_\phi(\mathcal{H} \mid \mathcal{A}_0^\ast)}_{\text{structure network}}
\cdot
\underbrace{q_\psi\big(\rmB \mid \rmY^\ast, \rmX^\ast, \Pi(\rmT), \mathcal{F}\big)}_{\text{CogFormer}},
\end{equation}
%
and marginalizing the structure yields a fully model-averaged parameter posterior,
%
\begin{equation}
\label{eq:td-bma}
q(\rmB \mid \mathcal{A}_0^\ast)
= \sum_{\mathcal{H} \in \mathbb{H}} q_\phi(\mathcal{H} \mid \mathcal{A}_0^\ast)\,
q_\psi(\rmB \mid \rmY^\ast, \rmX^\ast, \Pi(\rmT), \mathcal{F}),
\end{equation}
%
estimated without bias by drawing $\mathcal{H}^{(m)} \sim q_\phi$ and then
$\rmB^{(m)} \sim q_\psi(\cdot \mid \mathcal{H}^{(m)})$. \Cref{eq:td-bma} is Bayesian
model averaging over a space of $10^{10}$ models at the cost of $M$ forward passes.
Its practical content is that parameter uncertainty then includes structural
uncertainty, rather than conditioning on a structure that was itself selected from the
same data.

\paragraph{Decision rules.} Where a single theory must be reported, three summaries are
available: the MAP structure $\arg\max_\mathcal{H} q_\phi$; the \emph{median probability
structure} $\hat{T}_{c,j} = \mathbb{1}[\pi_{c,j} > 1/2]$, which is optimal for
prediction under squared-error loss in the classical linear setting and is the natural
default when marginal claims are the target; and the top-$k$ credible set, the smallest
set of structures whose posterior mass exceeds $1-\alpha$. The third is the most honest
and should be the default report: a broad posterior over theories \emph{is} the
scientific result when the data do not discriminate, and collapsing it to an argmax
manufactures a certainty the data do not contain.

%==============================================================================
\section{Traversing the structure space}
\label{sec:td-search}
%==============================================================================

The combinatorial size of $\mathbb{H}$ suggests that theory discovery is a search
problem. Under amortization it largely is not. Search algorithms exist because scoring a
hypothesis is expensive, so one must be economical about which hypotheses get scored.
Amortization removes that premise, and with it the algorithm class.

\subsection{Why the cost factorizes}

The decisive structural property is where the data-dependent work sits. Evaluating
$q_\phi(\mathcal{H} \mid \mathcal{A}_0^\ast)$ decomposes into
%
\begin{equation}
\label{eq:td-factorization}
\underbrace{\rmS = \mathrm{Enc}(\rmY^\ast, \rmX^\ast)}_{\mathcal{O}(N L d),\;\text{paid once, shared}}
\;\longrightarrow\;
\underbrace{q_\phi(\mathcal{H} \mid \rmS)}_{\mathcal{O}(n d^2),\; n = (I^\ast+1)D_\mathcal{F},\;\text{per hypothesis}} ,
\end{equation}
%
where $\rmS \in \mathbb{R}^{L \times S}$ are the encoder summaries. The expensive term
scales with the number of trials $N$ and is \emph{independent of the hypothesis}: it is
computed once and cached, then shared across every structure ever queried for this
dataset. Only the cheap term, which touches $n \approx 100$ decoder tokens and never
re-reads the data, is repeated.

Iterative methods admit no such factorization. Every refit of every node of the model
lattice re-reads all $N$ trials, so cost grows as (number of hypotheses) $\times$
(dataset size). Under \cref{eq:td-factorization} the two factors are additive rather
than multiplicative. This, rather than raw throughput, is why structure inference is
feasible at all: it is the same amortization argument that motivates CogFormer, applied
one level up.

\subsection{Sampling replaces search}

Because $q_\phi$ is a \emph{generative} model over $\mathbb{H}$, structures are drawn
ancestrally rather than sought. Under the any-order factorization
\cref{eq:td-aoar}, drawing $M$ structures requires $n$ \emph{sequential} decoder steps
at batch size $M$; latency is therefore $\mathcal{O}(n)$ and independent of $M$ up to
memory. Drawing $10^3$ theories costs what drawing one costs. The samples are
distributed as the posterior by construction, so the region of $\mathbb{H}$ that carries
mass is visited without any explicit traversal of the remaining $10^{10}$ hypotheses.

\subsection{Enumeration and ranking, when needed}

Explicit hypothesis lists are still wanted for reporting. Three regimes:

\begin{description}
\item[Factorized head: enumeration is closed-form.] Under \cref{eq:td-meanfield} the
joint factorizes, so
$\log q_\phi(\rmT) = \sum_{c,j} \log \pi_{c,j}^{T_{c,j}}(1-\pi_{c,j})^{1-T_{c,j}}$.
The MAP is read off cell-wise as $\mathbb{1}[\pi_{c,j} > 1/2]$, and the top-$k$
structures are obtained by flipping the $k-1$ cells with smallest $|\mathrm{logit}\,\pi_{c,j}|$
in ascending order of cost. No search is performed: the ranking is a sort of $n$ numbers.

\item[Autoregressive head: beam search.] Under \cref{eq:td-aoar} the MAP is a maximization
over a binary tree of depth $n$. Beam search with width $W$ recovers it approximately,
and exactly when the beam retains the optimum. Because beams batch, the cost is $n$
sequential decoder steps at batch $W$ --- the same $\mathcal{O}(n)$ latency as sampling.
Fixing the visitation order $\sigma$ to descending marginal certainty makes early
commitments the reliable ones.

\item[Restricted sub-lattices: exhaustive.] When a researcher restricts attention to a
handful of blocks and parameters, $|\mathbb{T}|$ falls to $2^{10}$--$2^{20}$ and the
autoregressive head can score the entire sub-lattice exactly in a single batched pass.
This is the regime of a conventional model-comparison table, recovered as a special
case.
\end{description}

\subsection{Propose and rerank}

The factorized and autoregressive heads compose into a two-stage scheme that inherits
the strengths of both. Draw $K$ candidates $\rmT^{(k)} \sim q_{\mathrm{MF}}$ from the
cheap factorized head, score each exactly under $q_{\mathrm{AR}}$, and form
self-normalized importance weights
%
\begin{equation}
\label{eq:td-snis}
w^{(k)} \propto \frac{q_{\mathrm{AR}}(\rmT^{(k)} \mid \mathcal{A}_0^\ast)}
                     {q_{\mathrm{MF}}(\rmT^{(k)} \mid \mathcal{A}_0^\ast)},
\qquad
\mathrm{ESS} = \Big(\textstyle\sum_k (w^{(k)})^2\Big)^{-1} .
\end{equation}
%
The weights correct the factorized proposal toward the joint, and the effective sample
size is a free diagnostic: $\mathrm{ESS} \approx K$ certifies that the structure
posterior is nearly separable and the PIPs may be read at face value, whereas
$\mathrm{ESS} \ll K$ localizes the dependence --- typically collinear blocks competing
for the same parameter --- and flags exactly the cases in which marginal inclusion
probabilities mislead.

\subsection{A zero-retraining route: encompassing-model screening}

The existing checkpoint already supports an approximate structure ranking, which is
useful as a baseline and as an immediate diagnostic. Query CogFormer with the full mask
$\rmM = \bm{1}$, the encompassing model in which every block is active. For nested
comparisons the Savage--Dickey density ratio gives the Bayes factor against inclusion
as a ratio of posterior to prior density evaluated at the origin,
%
\begin{equation}
\label{eq:td-savagedickey}
\mathrm{BF}_{0:1}(c,j)
= \frac{q_\psi\big(\beta_{c,j} = 0 \mid \rmY^\ast, \rmX^\ast, \bm{1}, \mathcal{F}\big)}
       {p\big(\beta_{c,j} = 0 \mid \bm{1}, \mathcal{F}\big)} ,
\end{equation}
%
requiring no fit of the reduced model. The denominator is analytic under the
variance-budget prior. The numerator is available because flow matching defines a
continuous normalizing flow: the log-density follows from the instantaneous change of
variables, $\log q_\psi(\rmB) = \log \mathcal{N}(\rmZ_1; \bm{0}, \mathbf{I}) + \int_0^1 \nabla \!\cdot\! \hat{u}\, dt$,
with the divergence estimated stochastically. Ranking cells by
\cref{eq:td-savagedickey} yields an ordered nested path of $n$ models rather than
$2^{n}$, which can then be scored individually.

Two caveats bound this route, and both are informative.

\begin{enumerate}
\item \emph{Prior compatibility.} Savage--Dickey requires
$p(\rmB_{-(c,j)} \mid \beta_{c,j} = 0, \bm{1}) = p(\rmB_{-(c,j)} \mid \rmT_{\mathrm{reduced}})$.
The variance-budget prior violates this: the Dirichlet split redistributes
$\sigma^2_{\beta,j}$ across whichever blocks are active, so removing a block changes the
prior on the survivors. The ratio is therefore approximate unless a nesting-compatible
variant of the coefficient prior is used.
\item \emph{Train--test mismatch.} During training the conditioning mask always equals
the generating mask, so the network never observes data generated under a sparse
structure but queried under a superset. The encompassing-model posterior is thus
partially extrapolated, and it is weakest precisely near $\beta = 0$, where
\cref{eq:td-savagedickey} is evaluated.
\end{enumerate}

The second caveat has a cheap fix that also points at the preferred architecture:
augment training by occasionally conditioning on a mask that strictly contains the
generating mask, with the surplus coefficients set to exactly zero. This places training
mass at the origin and makes the encompassing query in-distribution.

\subsection{The unified alternative: a gated flow}

Carrying the augmentation above to its conclusion replaces the two-network design with
one. Attach to each decoder token a Bernoulli inclusion gate alongside its velocity
output, so that the per-cell posterior is a spike-and-slab mixture
%
\begin{equation}
\label{eq:td-gatedflow}
q\big(\beta_{c,j} \mid \mathcal{A}_0\big)
= \big(1 - \pi_{c,j}\big)\,\delta_0(\beta_{c,j})
\;+\; \pi_{c,j}\, q_\psi\big(\beta_{c,j} \mid \rmT, \mathcal{A}_0\big),
\end{equation}
%
with the atom at zero carried by the gate---which a continuous flow cannot
represent---and the slab by the existing velocity field. Structure and coefficients are
then inferred by one network in one pass, the model-averaged posterior
\cref{eq:td-bma} becomes a single ancestral sampling sweep, and the two objectives share
an encoder rather than duplicating it.

\subsection{Cost summary}

\begin{table}[h]
\centering
\small
\begin{tabular}{@{}lll@{}}
\toprule
\textbf{Query} & \textbf{Sequential cost} & \textbf{Notes} \\
\midrule
Encoder summaries $\rmS$        & $\mathcal{O}(NLd)$, once & cached across all hypotheses \\
All PIPs (factorized)           & $1$ decoder pass         & marginal inclusion, closed-form top-$k$ \\
Score one enumerated $\mathcal{H}$ & $n$ decoder steps     & exact joint under \cref{eq:td-aoar} \\
Draw $M$ posterior structures   & $n$ steps, batch $M$     & latency independent of $M$ \\
MAP / top-$k$ (beam width $W$)  & $n$ steps, batch $W$     & latency independent of $W$ \\
Model-averaged $q(\rmB\mid\cdot)$ & $n$ steps $+$ ODE solve & \cref{eq:td-bma} \\
\midrule
Iterative refit of the lattice  & $\mathcal{O}(|\mathbb{H}| \cdot N \cdot \text{iters})$ & multiplicative in both factors \\
\bottomrule
\end{tabular}
\caption{Cost of structure queries, with $n = (I^\ast+1)D_\mathcal{F}$ the number of
theory cells. All amortized entries share one encoder evaluation.}
\label{tab:td-cost}
\end{table}

%==============================================================================
\section{The closed-world limit and the theory gap}
%==============================================================================

\Cref{eq:td-target} conditions on $\mathcal{H} \in \mathbb{H}$. The method therefore
returns the best theory \emph{in the specified space}, and its posterior is normalized
over that space regardless of whether any member fits. On real data this is the
dominant failure mode, and it is silent: a confident posterior over a badly
misspecified space looks exactly like a confident posterior over a well-specified one.
Two statistics bound the claim.

\paragraph{Closed-world adequacy.} Let $\rvs(\cdot)$ denote the encoder summary map and
$\tilde{\rmY}^{(m)}$ posterior-predictive replicates drawn through \cref{eq:td-bma}. The
discrepancy
%
\begin{equation}
\label{eq:td-mmd}
\Delta(\rmY^\ast) = \mathrm{MMD}^2\Big(\rvs(\rmY^\ast),\; \big\{\rvs(\tilde{\rmY}^{(m)})\big\}_{m=1}^M\Big)
\end{equation}
%
is calibrated against its distribution under simulated data from the meta-simulator,
yielding a null against which $\Delta(\rmY^\ast)$ can be tested. Rejection means no
theory in $\mathbb{H}$ reproduces the data, and the posterior over $\mathbb{H}$ should
not be interpreted.

\paragraph{Open-world theory gap.} Adequacy relative to the space is weaker than
adequacy in absolute terms. Let $\ell_\star$ be the held-out log predictive score of an
unrestricted black-box predictor of behavior---a forward meta-model such as Centaur,
which upper-bounds the predictable structure in the data---and $\ell_{\hat{\mathcal{H}}}$
the corresponding score of the selected theory. The \emph{theory gap}
%
\begin{equation}
\label{eq:td-gap}
\mathcal{R}(\rmY^\ast) = \ell_\star - \ell_{\hat{\mathcal{H}}}
\end{equation}
%
quantifies predictable structure that the best available interpretable theory fails to
capture. $\mathcal{R} \approx 0$ licenses the claim that $\hat{\mathcal{H}}$ is not
merely the best theory in $\mathbb{H}$ but adequate in absolute terms;
$\mathcal{R} \gg 0$ is positive evidence that a theory outside $\mathbb{H}$ exists, and
localizing $\mathcal{R}$ across trials or conditions indicates where.

%==============================================================================
\section{Validation}
%==============================================================================

Amortized posteriors over discrete structures require their own calibration checks;
parameter-space diagnostics do not transfer.

\begin{enumerate}
\item \textbf{Calibration of inclusion probabilities.} Bin simulated datasets by
predicted PIP and compare the bin mean to the empirical inclusion frequency. A
well-calibrated $q_\phi$ produces a diagonal reliability curve; this is the discrete
analogue of simulation-based calibration and is cheap to compute.
\item \textbf{Structure recovery.} Report Hamming distance between recovered and true
structures, and posterior rank of the true structure under \cref{eq:td-aoar}, as
functions of $N$, effect size (the design share $w_j$), and collinearity ($\zeta$,
$\eta^2$). The interesting result is the boundary: the regime in which the data cannot
identify the structure regardless of method.
\item \textbf{Reference agreement.} On families with tractable likelihoods, compare
\cref{eq:td-bayesfactor} against bridge-sampled Bayes factors on small enumerable
sub-lattices. This is the only direct check that the learned evidence ratios are right.
\item \textbf{Consensus recovery on real data.} On benchmark datasets where the field
agrees on the qualitative conclusion---speed--accuracy instructions on threshold,
stimulus difficulty on drift rate---the posterior should recover the consensus without
being told the design configuration.
\end{enumerate}

%==============================================================================
\section{Extension: designing the discriminating experiment}
%==============================================================================

Because the design matrix is a conditioning variable, the same machinery scores
\emph{prospective} designs. The expected information gain about the theory from a
candidate design $\rmX'$,
%
\begin{equation}
\label{eq:td-eig}
\mathrm{EIG}(\rmX') = \mathbb{E}_{p(\rmY' \mid \rmX', \mathcal{A}_0^\ast)}
\Big[ \mathbb{H}\big[\mathcal{H} \mid \mathcal{A}_0^\ast\big]
- \mathbb{H}\big[\mathcal{H} \mid \mathcal{A}_0^\ast, \rmY', \rmX'\big] \Big],
\end{equation}
%
is estimable by simulating $\rmY'$ from the current theory posterior and re-evaluating
$q_\phi$, which requires only forward passes. Maximizing \cref{eq:td-eig} selects the
experiment that best discriminates among the theories still in contention---closing the
loop from data to theory to next experiment.

\end{document}
